\nonstopmode{}
\documentclass[a4paper]{book}
\usepackage[times,inconsolata,hyper]{Rd}
\usepackage{makeidx}
\makeatletter\@ifl@t@r\fmtversion{2018/04/01}{}{\usepackage[utf8]{inputenc}}\makeatother
% \usepackage{graphicx} % @USE GRAPHICX@
\makeindex{}
\begin{document}
\chapter*{}
\begin{center}
{\textbf{\huge Package `sleepcleanr'}}
\par\bigskip{\large \today}
\end{center}
\ifthenelse{\boolean{Rd@use@hyper}}{\hypersetup{pdftitle = {sleepcleanr: Sleep EMA Diary Data Cleaning Pipeline}}}{}
\ifthenelse{\boolean{Rd@use@hyper}}{\hypersetup{pdfauthor = {Cai Dong; Maia ten Brink}}}{}
\begin{description}
\raggedright{}
\item[Title]\AsIs{Sleep EMA Diary Data Cleaning Pipeline}
\item[Version]\AsIs{1.4.5}
\item[Author]\AsIs{Cai Dong [aut, cre], Maia ten Brink [aut]}
\item[Maintainer]\AsIs{Cai Dong }\email{cyracaid@gmail.com}\AsIs{}
\item[Description]\AsIs{Reproducible R pipeline for cleaning sleep EMA diary data:
parses raw timestamps, detects and corrects temporal and duration errors
through a documented human-in-the-loop workflow, computes standard sleep
metrics (TST, SOL, WASO, SE), and generates diagnostic and research-ready
figures. Fully configurable via YAML so new studies can map their data
without code changes. Developed for the Stanford Psychophysiology
Laboratory's intensive-longitudinal sleep study (github.com/stanford-sber),
which supplied the data requirements.}
\item[License]\AsIs{MIT + file LICENSE}
\item[Encoding]\AsIs{UTF-8}
\item[URL]\AsIs{}\url{https://github.com/cyracaid/sleepdiary-cleaner}\AsIs{,
}\url{https://homesleep.stanford.edu}\AsIs{,
}\url{https://github.com/stanford-sber}\AsIs{}
\item[BugReports]\AsIs{}\url{https://github.com/cyracaid/sleepdiary-cleaner/issues}\AsIs{}
\item[Depends]\AsIs{R (>= 4.2)}
\item[Imports]\AsIs{corrplot, dplyr, ggplot2, gtable, lubridate, magick,
patchwork, readr, rlang, stringr, tidyr, utils, yaml}
\item[Suggests]\AsIs{gridExtra, knitr, rmarkdown, testthat (>= 3.0.0), withr}
\item[VignetteBuilder]\AsIs{knitr}
\item[Config/testthat/edition]\AsIs{3}
\item[Config/roxygen2/version]\AsIs{8.1.0}
\end{description}
\Rdcontents{Contents}
\HeaderA{.step\_ledger\_env}{Per-step flag ledger (log\_step)}{.step.Rul.ledger.Rul.env}
%
\begin{Description}
Intercepts EVERY pipeline step and records, against the shared final
standards (see flag\_standards.R), how many records fall in each flag category
at that step. Produces the data behind the new Figure 12 step x flag table.
\end{Description}
%
\begin{Usage}
\begin{verbatim}
.step_ledger_env
\end{verbatim}
\end{Usage}
%
\begin{Details}
Two kinds of reduction are tracked separately:
- n\_corrected : rows fixed by a manual/auto correction (data changed)
- n\_suppressed: rows a human accepted as not-an-error (label withdrawn,
data unchanged)
\end{Details}
\HeaderA{adapt\_columns}{Apply column mapping to a data frame}{adapt.Rul.columns}
%
\begin{Description}
Renames columns in \code{data} according to the mapping defined in config.
Columns whose mapped name is NULL are skipped.
\end{Description}
%
\begin{Usage}
\begin{verbatim}
adapt_columns(data, config)
\end{verbatim}
\end{Usage}
%
\begin{Arguments}
\begin{ldescription}
\item[\code{data}] Data frame. Raw input data with user's column names.

\item[\code{config}] List. Configuration list from \code{load\_config()}.
\end{ldescription}
\end{Arguments}
%
\begin{Value}
Data frame with columns renamed to pipeline-internal names.
\end{Value}
\HeaderA{as.data.frame.sleep\_diary}{Extract the working data frame from a sleep\_diary}{as.data.frame.sleep.Rul.diary}
%
\begin{Description}
The escape hatch for backward compatibility: any v1.2.0 code that expects a
plain data frame can call this and carry on unchanged.
\end{Description}
%
\begin{Usage}
\begin{verbatim}
## S3 method for class 'sleep_diary'
as.data.frame(x, row.names = NULL, optional = FALSE, ...)
\end{verbatim}
\end{Usage}
%
\begin{Arguments}
\begin{ldescription}
\item[\code{x}] A \code{sleep\_diary} object.

\item[\code{row.names}] NULL or a character vector giving row names.

\item[\code{optional}] Logical. Unused; present for S3 generic consistency.

\item[\code{...}] Unused.
\end{ldescription}
\end{Arguments}
%
\begin{Value}
A data frame.
\end{Value}
\HeaderA{assert\_contract\_columns}{Assert that a sleep\_diary carries the public contract columns}{assert.Rul.contract.Rul.columns}
%
\begin{Description}
The interface contract Step 7 promises downstream consumers. Kept separate
from \code{step\_compute\_metrics()} so tests can assert the contract without
re-running the step.
\end{Description}
%
\begin{Usage}
\begin{verbatim}
assert_contract_columns(x, error = TRUE)
\end{verbatim}
\end{Usage}
%
\begin{Arguments}
\begin{ldescription}
\item[\code{x}] A \code{sleep\_diary} object or data frame.

\item[\code{error}] Logical. If TRUE (default) raise on a missing column; if FALSE
return the missing names.
\end{ldescription}
\end{Arguments}
%
\begin{Value}
Character vector of missing columns, invisibly.
\end{Value}
\HeaderA{as\_sleep\_diary}{Coerce a data frame into a sleep\_diary}{as.Rul.sleep.Rul.diary}
%
\begin{Description}
Coerce a data frame into a sleep\_diary
\end{Description}
%
\begin{Usage}
\begin{verbatim}
as_sleep_diary(x, ...)
\end{verbatim}
\end{Usage}
%
\begin{Arguments}
\begin{ldescription}
\item[\code{x}] A data frame, or an existing sleep\_diary (returned unchanged).

\item[\code{...}] Passed to \code{new\_sleep\_diary()}.
\end{ldescription}
\end{Arguments}
%
\begin{Value}
A \code{sleep\_diary} object.
\end{Value}
\HeaderA{bland\_altman}{Bland-Altman analysis and threshold validation}{bland.Rul.altman}
%
\begin{Description}
\code{bland\_altman()} compares two measurement methods using the
Bland-Altman (1986) method.  It is advisory only: it does not modify
data and is never fed back into the pipeline decision tree.

\code{\LinkA{validate\_thresholds}{validate.Rul.thresholds}} compares every configured cleaning
threshold against the 95\% limits-of-agreement half-width to assess
whether each cutoff sits safely above measurement noise.
\end{Description}
%
\begin{Usage}
\begin{verbatim}
bland_altman(data, reported_col, computed_col,
             label = "Metric", loa_ci = FALSE)
\end{verbatim}
\end{Usage}
%
\begin{Arguments}
\begin{ldescription}
\item[\code{data}] A data frame (typically after Step 7).
\item[\code{reported\_col}] Name of the self-reported column.
\item[\code{computed\_col}] Name of the pipeline-computed column.
\item[\code{label}] Human-readable metric name for plot titles.
\item[\code{loa\_ci}] Logical. Compute 95\% CI for limits of agreement.
\end{ldescription}
\end{Arguments}
%
\begin{Value}
\code{bland\_altman()} returns a list with components \code{bias},
\code{lower\_loa}, \code{upper\_loa}, \code{prop\_bias\_p}, \code{n\_pairs}
and \code{plot}.
\end{Value}
%
\begin{References}
Bland, J. M. and Altman, D. G. (1986). Statistical methods for
assessing agreement between two methods of clinical measurement.
\emph{The Lancet}, 327(8476), 307-310.
\end{References}
\HeaderA{cfg\_get}{Safe config\_get -- fetches pipeline\_config from global env automatically}{cfg.Rul.get}
%
\begin{Description}
Use this in standalone scripts (sleep\_visualization.R, checkforerrors\_processing.R)
where pipeline\_config may not exist in the calling scope.
\end{Description}
%
\begin{Usage}
\begin{verbatim}
cfg_get(key, default = NULL, cfg = NULL)
\end{verbatim}
\end{Usage}
%
\begin{Arguments}
\begin{ldescription}
\item[\code{key}] Character. Dot-separated key.

\item[\code{default}] Default value if key not found.

\item[\code{cfg}] Optional. A pipeline configuration list. When provided, this
is used directly instead of falling back to the global environment.
**From v1.3.1, passing \code{cfg} explicitly is the preferred path.**
The global-environment fallback is deprecated and will emit a warning.
\end{ldescription}
\end{Arguments}
%
\begin{Value}
The config value, or \code{default}.
\end{Value}
\HeaderA{cleaning\_chain}{Pipeline step adapters and the cleaning chain}{cleaning.Rul.chain}
%
\begin{Description}
The step adapters wrap the v1.2.0 pipeline scripts without changing
their logic, adding timing, column-diff tracking and ledger logging.
\code{run\_cleaning\_chain()} composes them into a single callable chain.
\end{Description}
%
\begin{Details}
Each adapter is documented individually: \code{\LinkA{run\_cleaning\_chain}{run.Rul.cleaning.Rul.chain}},
\code{\LinkA{assert\_contract\_columns}{assert.Rul.contract.Rul.columns}},
\code{\LinkA{step\_process\_timestamps}{step.Rul.process.Rul.timestamps}}, \code{\LinkA{step\_process\_intervals}{step.Rul.process.Rul.intervals}},
\code{\LinkA{step\_normalize\_sequence}{step.Rul.normalize.Rul.sequence}}, \code{\LinkA{step\_apply\_corrections}{step.Rul.apply.Rul.corrections}},
\code{\LinkA{step\_apply\_duration\_corrections}{step.Rul.apply.Rul.duration.Rul.corrections}},
\code{\LinkA{step\_compute\_metrics}{step.Rul.compute.Rul.metrics}}.
\end{Details}
\HeaderA{config\_col}{Get column mapping from config}{config.Rul.col}
\keyword{internal}{config\_col}
%
\begin{Description}
Returns the user's column name for a given pipeline-internal column.
\end{Description}
%
\begin{Usage}
\begin{verbatim}
config_col(config, internal_name)
\end{verbatim}
\end{Usage}
%
\begin{Arguments}
\begin{ldescription}
\item[\code{config}] List. Configuration list.

\item[\code{internal\_name}] Character. Pipeline-internal column name.
\end{ldescription}
\end{Arguments}
%
\begin{Value}
Character. User's column name, or \code{internal\_name} if not mapped.
\end{Value}
\HeaderA{config\_get}{Get a nested config value by dot-separated key}{config.Rul.get}
\keyword{internal}{config\_get}
%
\begin{Description}
Get a nested config value by dot-separated key
\end{Description}
%
\begin{Usage}
\begin{verbatim}
config_get(config, key, default = NULL)
\end{verbatim}
\end{Usage}
%
\begin{Arguments}
\begin{ldescription}
\item[\code{config}] List. Configuration list from \code{load\_config()}.

\item[\code{key}] Character. Dot-separated key, e.g. \code{"classification.temporal.max\_sol\_minutes"}.

\item[\code{default}] Default value if key not found.
\end{ldescription}
\end{Arguments}
%
\begin{Value}
The config value, or \code{default}.
\end{Value}
\HeaderA{dim.sleep\_diary}{Dimensions of a sleep\_diary}{dim.sleep.Rul.diary}
%
\begin{Description}
Defined so that \code{nrow()} and \code{ncol()} work directly on the object.
Note that \code{nrow} itself is not generic in base R -- it dispatches
through \code{dim()}, which is why this method exists rather than an
\code{nrow.sleep\_diary}.
\end{Description}
%
\begin{Usage}
\begin{verbatim}
## S3 method for class 'sleep_diary'
dim(x)
\end{verbatim}
\end{Usage}
%
\begin{Arguments}
\begin{ldescription}
\item[\code{x}] A \code{sleep\_diary} object.
\end{ldescription}
\end{Arguments}
%
\begin{Value}
Integer vector of length 2: rows, columns.
\end{Value}
\HeaderA{evaluate\_all\_standards}{Evaluate ALL standards, returning per-record label data frame.}{evaluate.Rul.all.Rul.standards}
%
\begin{Description}
Evaluate ALL standards, returning per-record label data frame.
\end{Description}
%
\begin{Usage}
\begin{verbatim}
evaluate_all_standards(df, cfg = NULL)
\end{verbatim}
\end{Usage}
%
\begin{Arguments}
\begin{ldescription}
\item[\code{df}] A data frame to evaluate all standards against.

\item[\code{cfg}] Pipeline configuration list.
\end{ldescription}
\end{Arguments}
\HeaderA{eval\_checkforerrors}{Evaluate checkforerrors (auto-detection flags)}{eval.Rul.checkforerrors}
%
\begin{Description}
Prereq: Step 8. TIMESTAMP\_ISSUE | DURATION\_ISSUE | AMOUNT\_FLAG |
SELF\_REPORTED\_FLAG | CLEAN | NEEDS\_REVIEW.
\end{Description}
%
\begin{Usage}
\begin{verbatim}
eval_checkforerrors(df, cfg = NULL)
\end{verbatim}
\end{Usage}
%
\begin{Arguments}
\begin{ldescription}
\item[\code{df}] A data frame with checkforerrors columns.

\item[\code{cfg}] Pipeline configuration list.
\end{ldescription}
\end{Arguments}
\HeaderA{eval\_data\_category}{Evaluate data\_category (temporal classification)}{eval.Rul.data.Rul.category}
%
\begin{Description}
Authoritative from Step 5. Categories: clean | error | unusual | equal\_time\_ok |
reasonable\_unusual | skipped\_na.
\end{Description}
%
\begin{Usage}
\begin{verbatim}
eval_data_category(df, cfg = NULL)
\end{verbatim}
\end{Usage}
%
\begin{Arguments}
\begin{ldescription}
\item[\code{df}] A data frame with corrected timestamp columns.

\item[\code{cfg}] Pipeline configuration list.
\end{ldescription}
\end{Arguments}
\HeaderA{eval\_duration\_extreme}{Evaluate duration\_extreme (separate from severity count)}{eval.Rul.duration.Rul.extreme}
%
\begin{Description}
Evaluate duration\_extreme (separate from severity count)
\end{Description}
%
\begin{Usage}
\begin{verbatim}
eval_duration_extreme(df, cfg = NULL)
\end{verbatim}
\end{Usage}
%
\begin{Arguments}
\begin{ldescription}
\item[\code{df}] A data frame with sleep\_duration\_h column.

\item[\code{cfg}] Pipeline configuration list.
\end{ldescription}
\end{Arguments}
\HeaderA{eval\_field\_misentry}{Evaluate field\_misentry (cross-participant field misentry, Step 1.5)}{eval.Rul.field.Rul.misentry}
%
\begin{Description}
Evaluate field\_misentry (cross-participant field misentry, Step 1.5)
\end{Description}
%
\begin{Usage}
\begin{verbatim}
eval_field_misentry(df, cfg = NULL)
\end{verbatim}
\end{Usage}
%
\begin{Arguments}
\begin{ldescription}
\item[\code{df}] A data frame with SOL/WASO and timestamp columns.

\item[\code{cfg}] Pipeline configuration list.
\end{ldescription}
\end{Arguments}
\HeaderA{eval\_flag\_severity}{Evaluate flag\_severity (computed metric flags)}{eval.Rul.flag.Rul.severity}
%
\begin{Description}
Prereq: Step 7 metrics. Clean (0) | Minor issues (1 flag) | Major issues (2+ flags).
\end{Description}
%
\begin{Usage}
\begin{verbatim}
eval_flag_severity(df, cfg = NULL)
\end{verbatim}
\end{Usage}
%
\begin{Arguments}
\begin{ldescription}
\item[\code{df}] A data frame with sleep\_efficiency\_pct, sol\_h, waso\_h columns.

\item[\code{cfg}] Pipeline configuration list.
\end{ldescription}
\end{Arguments}
\HeaderA{figure12\_step\_flag\_table}{Figure 12 (new) --- Step x Flag ledger table}{figure12.Rul.step.Rul.flag.Rul.table}
%
\begin{Description}
Replaces the coarse A-E bar chart. Renders one row per pipeline step and,
against the shared final standards, shows how each flag family is generated
and reduced across steps. "not computable at this step" shows as "---".
\end{Description}
%
\begin{Usage}
\begin{verbatim}
figure12_step_flag_table(
  cfg = NULL,
  output_dir = ".",
  save_png = NULL,
  filename = "12_Pipeline_Correction_Progress"
)
\end{verbatim}
\end{Usage}
%
\begin{Arguments}
\begin{ldescription}
\item[\code{cfg}] Pipeline configuration list.

\item[\code{output\_dir}] Directory for saving output PNG.

\item[\code{save\_png}] Optional save function for PNG output.

\item[\code{filename}] Output PNG filename without extension.
\end{ldescription}
\end{Arguments}
%
\begin{Details}
Drop-in: replace the current Figure 12 block in sleep\_visualization.R with a
call to `figure12\_step\_flag\_table(cfg = cfg, output\_dir = output\_dir, save\_png = save\_png)`.
\end{Details}
\HeaderA{figure\_cleaning\_effect}{Figure 2 --- Effect of cleaning (before vs after)}{figure.Rul.cleaning.Rul.effect}
%
\begin{Description}
Three-panel figure showing what cleaning actually changed.
\end{Description}
%
\begin{Usage}
\begin{verbatim}
figure_cleaning_effect()
\end{verbatim}
\end{Usage}
%
\begin{Details}
The three panels are:
\begin{description}

\item[A] SOL before (self-reported) versus after (computed), drawn as a
combined boxplot and violin.
\item[B] TST distribution after cleaning, stacked by flag severity.
\item[C] Individual record changes --- self-reported versus computed SOL,
coloured by correction type (automatic, manual, unchanged). This panel
is the one that shows only corrected observations moved.

\end{description}


Takes no arguments. All data is extracted from \code{corrected\_ema\_data} in
the global environment, so the pipeline must have been run first ---
otherwise the function stops with an explanatory message.

Requires the \pkg{patchwork} package for the combined layout.
\end{Details}
%
\begin{Value}
The combined \code{ggplot} object, invisibly. Called for its side effect of
writing \file{figures/Figure\_2\_Cleaning\_Effect.png}.
\end{Value}
%
\begin{SeeAlso}
\code{\LinkA{figure\_pipeline\_workflow}{figure.Rul.pipeline.Rul.workflow}} for the companion flow diagram, and
\code{\LinkA{run\_pipeline}{run.Rul.pipeline}} to produce the required input.
\end{SeeAlso}
\HeaderA{figure\_pipeline\_workflow}{Figure 1 --- Pipeline workflow flow diagram}{figure.Rul.pipeline.Rul.workflow}
%
\begin{Description}
Produces a publication-quality left-to-right flow diagram showing how raw
records move through automatic validation, algorithmic correction, manual
review, and into the final clean dataset. Every box carries a record count
and its percentage of the raw total.
\end{Description}
%
\begin{Usage}
\begin{verbatim}
figure_pipeline_workflow()
\end{verbatim}
\end{Usage}
%
\begin{Details}
Takes no arguments. All counts are extracted from \code{corrected\_ema\_data}
in the global environment, so the pipeline must have been run first ---
otherwise the function stops with an explanatory message.

The figure is written to \file{figures/Figure\_1\_Pipeline\_Workflow.png} at
11 x 6.5 inches, 300 dpi, suitable for a manuscript Methods section.

Note that counts are reconstructed from the flag columns of the final
dataset, so they describe records that survived to the end of the pipeline.
\end{Details}
%
\begin{Value}
The \code{ggplot} object, invisibly. Called for its side effect of
writing \file{figures/Figure\_1\_Pipeline\_Workflow.png}.
\end{Value}
%
\begin{SeeAlso}
\code{\LinkA{figure\_cleaning\_effect}{figure.Rul.cleaning.Rul.effect}} for the companion before/after figure,
and \code{\LinkA{run\_pipeline}{run.Rul.pipeline}} to produce the required input.
\end{SeeAlso}
\HeaderA{figure\_run\_dir}{Resolve the figure output directory for a run}{figure.Rul.run.Rul.dir}
%
\begin{Description}
Builds the run-specific visualization directory name:
\code{<root>/latest\_visualization\_<tag>\_n<records>}. The base root comes from
config (\code{output.figure.root\_dir} for real data, \code{output.figure.synth\_root\_dir}
for synthetic/test AND unresolved data). Only a data\_tag that has been
positively confirmed as \code{"real"} is trusted with the \code{output/}
root -- \code{"synth"} and \code{"unknown"} both stay outside \code{output/}
by default. This is a fail-closed choice: \code{"unknown"} means the tag
detector could not read \code{data.files.main} from the config (see
\code{sleep\_visualization.R}), i.e. we genuinely do not know what was
loaded. Routing it next to confirmed real-data figures would silently
extend real-data trust (gitignored, "never leaves the study team" handling)
to data whose identity is unverified. Before 2026-08-11, "unknown" fell
into the \code{else} branch and was treated exactly like "real"; that has
been fixed so only "real" gets the privileged path.
\end{Description}
%
\begin{Usage}
\begin{verbatim}
figure_run_dir(cfg = NULL, data_tag, n_records = NULL)
\end{verbatim}
\end{Usage}
%
\begin{Arguments}
\begin{ldescription}
\item[\code{cfg}] Optional. Pipeline configuration list (preferred; falls back to
the global environment when omitted).

\item[\code{data\_tag}] Character. "real", "synth", or "unknown".

\item[\code{n\_records}] Numeric or NULL. Row count appended to the directory name.
\end{ldescription}
\end{Arguments}
%
\begin{Value}
Character. Relative path to the figure output directory.
\end{Value}
\HeaderA{finalize\_columns}{Build the analysis-facing datasets from the full pipeline output}{finalize.Rul.columns}
%
\begin{Description}
Splits the pipeline output into two delivered datasets plus the full
archive, driven entirely by `inst/extdata/column\_dictionary.csv`.
\end{Description}
%
\begin{Usage}
\begin{verbatim}
finalize_columns(
  data,
  review_data = NULL,
  dict_path = NULL,
  output_dir = "output",
  write = TRUE,
  verbose = TRUE
)
\end{verbatim}
\end{Usage}
%
\begin{Arguments}
\begin{ldescription}
\item[\code{data}] The pipeline output (`corrected\_ema\_data`).

\item[\code{review\_data}] Optional. `review\_output\$data\_with\_flags`. Step 8 creates
its flags on a copy and never writes them back to `corrected\_ema\_data`
(see `00\_MAIN\_entry.R`, Step 8), so columns such as `needs\_review\_flag`
are only reachable through this object. The dictionary's `source\_object`
field records which columns need it.

\item[\code{dict\_path}] Path to the column dictionary CSV. Defaults to the copy
shipped in `inst/extdata`.

\item[\code{output\_dir}] Directory for the delivered files.

\item[\code{write}] Write files to disk. Set FALSE to inspect the return
value without touching the filesystem.

\item[\code{verbose}] Print a summary.
\end{ldescription}
\end{Arguments}
%
\begin{Details}
The dictionary is the single source of truth for three things that used to
drift apart: the column whitelist, the rename mapping, and the data
dictionary itself. Adding a column means editing one CSV row, not three
places.
\end{Details}
%
\begin{Value}
Invisibly, a list with `final` (Dataset A), `prepost` (Dataset B)
and `full` (everything, unchanged).
\end{Value}
\HeaderA{flag\_statistical\_outliers}{Per-participant statistical outlier detection via IQR}{flag.Rul.statistical.Rul.outliers}
%
\begin{Description}
Detects within-participant outliers on key sleep metrics using the
interquartile range (IQR) method: any value more than \code{1.5 * IQR}
below Q1 or above Q3 is flagged.
\end{Description}
%
\begin{Usage}
\begin{verbatim}
flag_statistical_outliers(
  data,
  group_col = "pid",
  metrics = c(self_diffcalc_sol_minutes = "sol", self_diffcalc_totalsleeptime_minutes =
    "tst", self_diffcalc_sleepefficiency_percent = "se", avg_waso_estimate_am_minutes =
    "waso"),
  multiplier = 1.5,
  min_rows = 5L
)
\end{verbatim}
\end{Usage}
%
\begin{Arguments}
\begin{ldescription}
\item[\code{data}] A data frame from the pipeline (typically after Step 7).

\item[\code{group\_col}] Character. Column identifying participants (default
\code{"pid"}).

\item[\code{metrics}] Character vector. Columns to scan for outliers.
Defaults to the four primary derived metrics.

\item[\code{multiplier}] Numeric. The IQR multiplier. Default 1.5 (Tukey's
standard). Use 3.0 for "far out" detection only.

\item[\code{min\_rows}] Integer. Participants with fewer than this many valid
rows are skipped (their IQR is unreliable).
\end{ldescription}
\end{Arguments}
%
\begin{Details}
This complements (does NOT replace) the hard-threshold system in the
cleaning pipeline. Hard thresholds detect values that are implausible
in absolute terms (e.g. SOL > 120 min). IQR flags detect values that
are unusual \emph{for that specific participant} -- a person who
typically falls asleep in 5 minutes suddenly reporting 45 minutes is
worth a closer look even though 45 min is below the hard cutoff.
\end{Details}
%
\begin{Value}
A copy of \code{data} with one new column:
\begin{description}

\item[\code{iqr\_outlier}] Character. Names of the metrics that were
out of range, separated by "+" (e.g. \code{"sol+tst"}). \code{NA}
when no metric is flagged or the participant has too few rows.

\end{description}

\end{Value}
%
\begin{Examples}
\begin{ExampleCode}
## Not run: 
flagged <- flag_statistical_outliers(corrected_ema_data)
table(flagged$iqr_outlier)

## End(Not run)

\end{ExampleCode}
\end{Examples}
\HeaderA{get\_step\_ledger\_long}{Flatten the ledger into a long data frame}{get.Rul.step.Rul.ledger.Rul.long}
%
\begin{Description}
Flatten the ledger into a long data frame
\end{Description}
%
\begin{Usage}
\begin{verbatim}
get_step_ledger_long()
\end{verbatim}
\end{Usage}
\HeaderA{get\_step\_ledger\_wide}{Wide ledger for one standard}{get.Rul.step.Rul.ledger.Rul.wide}
%
\begin{Description}
Wide ledger for one standard
\end{Description}
%
\begin{Usage}
\begin{verbatim}
get_step_ledger_wide(standard = "data_category")
\end{verbatim}
\end{Usage}
%
\begin{Arguments}
\begin{ldescription}
\item[\code{standard}] Character. The standard to pivot wide. Default: "data\_category".
\end{ldescription}
\end{Arguments}
\HeaderA{handle\_missing}{Tag missing-data reason codes and optionally carry forward single-day gaps}{handle.Rul.missing}
%
\begin{Description}
Adds a \code{missing\_reason} column that distinguishes why a row has
incomplete data, rather than lumping everything into \code{skipped\_na}.
Optionally applies last-observation-carried-forward (LOCF) to metric
columns for single-day gaps.
\end{Description}
%
\begin{Usage}
\begin{verbatim}
handle_missing(
  data,
  timestamp_cols = NULL,
  metric_cols = NULL,
  group_col = "pid",
  order_col = "day_num",
  max_gap = 1L
)
\end{verbatim}
\end{Usage}
%
\begin{Arguments}
\begin{ldescription}
\item[\code{data}] A data frame from the pipeline.

\item[\code{timestamp\_cols}] Character. Columns that carry the four sleep-event
POSIXct timestamps. If \code{NULL}, auto-detects from common patterns.

\item[\code{metric\_cols}] Character. Columns eligible for LOCF. If \code{NULL},
auto-detects the standard derived-metric columns.

\item[\code{group\_col}] Character. Participant identifier. Default \code{"pid"}.

\item[\code{order\_col}] Character. Within-participant ordering column. Default
\code{"day\_num"}.

\item[\code{max\_gap}] Integer. Maximum consecutive missing days to fill via
LOCF. Default 1. Set to 0 to disable LOCF entirely.
\end{ldescription}
\end{Arguments}
%
\begin{Value}
A copy of \code{data} with:
\begin{description}

\item[\code{missing\_reason}] Character. One of \code{"all\_timestamps\_na"},
\code{"partial\_timestamps\_na"}, \code{"derived\_na"}, or \code{NA}
(complete row).
\item[For each metric column when \code{max\_gap > 0}:] a companion
\code{<col>\_imputed} logical column, and the original column may
contain LOCF-filled values.

\end{description}

\end{Value}
%
\begin{Section}{Design boundaries}


* Timestamp columns are **never** imputed -- filling a missing bedtime
would fabricate an event that did not happen.
* Only metric columns (SOL, WASO, TST, SE) are eligible for LOCF.
* LOCF is capped at \code{max\_gap} consecutive days (default 1) to
avoid filling week-long gaps with stale data.
* All imputed values carry an \code{\_imputed} companion column
(\code{TRUE} / \code{FALSE}) for audit trail.
\end{Section}
%
\begin{Examples}
\begin{ExampleCode}
## Not run: 
handled <- handle_missing(corrected_ema_data)
table(handled$missing_reason)

# LOCF disabled -- only reason codes
handled <- handle_missing(corrected_ema_data, max_gap = 0)

## End(Not run)

\end{ExampleCode}
\end{Examples}
\HeaderA{init\_step\_ledger}{Start a fresh ledger (call once at the top of run\_pipeline).}{init.Rul.step.Rul.ledger}
%
\begin{Description}
Start a fresh ledger (call once at the top of run\_pipeline).
\end{Description}
%
\begin{Usage}
\begin{verbatim}
init_step_ledger()
\end{verbatim}
\end{Usage}
\HeaderA{is\_sleep\_diary}{Test whether an object is a sleep\_diary}{is.Rul.sleep.Rul.diary}
%
\begin{Description}
Test whether an object is a sleep\_diary
\end{Description}
%
\begin{Usage}
\begin{verbatim}
is_sleep_diary(x)
\end{verbatim}
\end{Usage}
%
\begin{Arguments}
\begin{ldescription}
\item[\code{x}] Any object.
\end{ldescription}
\end{Arguments}
%
\begin{Value}
Logical scalar.
\end{Value}
\HeaderA{load\_config}{Load pipeline configuration}{load.Rul.config}
%
\begin{Description}
Reads a YAML config file and returns a list of settings.
Falls back to the bundled default config if no file is specified.
\end{Description}
%
\begin{Usage}
\begin{verbatim}
load_config(config_file = NULL)
\end{verbatim}
\end{Usage}
%
\begin{Arguments}
\begin{ldescription}
\item[\code{config\_file}] Character. Path to a YAML config file, or NULL to use
the bundled default (\code{inst/config\_default.yaml}).
\end{ldescription}
\end{Arguments}
%
\begin{Value}
List of pipeline configuration values.
\end{Value}
\HeaderA{log\_step}{Record the flag state after a step.}{log.Rul.step}
%
\begin{Description}
Record the flag state after a step.
\end{Description}
%
\begin{Usage}
\begin{verbatim}
log_step(df, step_id, label, cfg = NULL, verbose = TRUE)
\end{verbatim}
\end{Usage}
%
\begin{Arguments}
\begin{ldescription}
\item[\code{df}] Data frame in its state AFTER the step.

\item[\code{step\_id}] Short ordered id, e.g. "1", "1.5", "2", ... "8.5".

\item[\code{label}] Human-readable step name.

\item[\code{cfg}] Config list (for thresholds).

\item[\code{verbose}] Logical. Print progress messages. Default: TRUE.
\end{ldescription}
\end{Arguments}
%
\begin{Value}
invisibly the row that was appended.
\end{Value}
\HeaderA{missing\_handler}{Missing-data reason codes and single-day LOCF}{missing.Rul.handler}
%
\begin{Description}
Adds a \code{missing\_reason} column that distinguishes why a row has
incomplete data.  Optionally applies last-observation-carried-forward
(LOCF) to metric columns for single-day gaps.  Timestamp columns are
never imputed; all filled values carry a \code{\_imputed} companion
column for audit trail.
\end{Description}
%
\begin{Details}
\code{\LinkA{handle\_missing}{handle.Rul.missing}} and \code{\LinkA{summarise\_missing}{summarise.Rul.missing}} are
documented individually.
\end{Details}
\HeaderA{new\_sleep\_diary}{Construct a sleep\_diary object}{new.Rul.sleep.Rul.diary}
%
\begin{Description}
Construct a sleep\_diary object
\end{Description}
%
\begin{Usage}
\begin{verbatim}
new_sleep_diary(
  data,
  step_id = "0",
  step_label = "raw",
  cfg = NULL,
  history = list(),
  extra = list()
)
\end{verbatim}
\end{Usage}
%
\begin{Arguments}
\begin{ldescription}
\item[\code{data}] A data frame holding the working records.

\item[\code{step\_id}] Character. Short ordered step id, e.g. "1", "1.5", "2".
Use "0" for a freshly loaded object that no step has processed yet.

\item[\code{step\_label}] Character. Human-readable step name.

\item[\code{cfg}] List or NULL. Pipeline configuration in force.

\item[\code{history}] List of prior step records (oldest first).

\item[\code{extra}] List. Optional extra fields merged into the step record,
for example \code{list(n\_corrected = 12)}.
\end{ldescription}
\end{Arguments}
%
\begin{Value}
An object of class \code{sleep\_diary}.
\end{Value}
\HeaderA{outlier\_flags}{Per-participant IQR outlier detection}{outlier.Rul.flags}
%
\begin{Description}
Detects within-participant outliers on key sleep metrics using the
interquartile range (IQR) method.  This complements (does not replace)
the hard-threshold system: hard thresholds catch implausible absolute
values, while IQR flags catch values that are unusual for that
specific participant.
\end{Description}
%
\begin{Details}
\code{\LinkA{flag\_statistical\_outliers}{flag.Rul.statistical.Rul.outliers}} and \code{\LinkA{summarise\_outliers}{summarise.Rul.outliers}}
are documented individually.
\end{Details}
\HeaderA{pipeline\_steps}{Pipeline step adapters (wrapper layer)}{pipeline.Rul.steps}
%
\begin{Description}
Every function here wraps ONE existing pipeline script from \code{v1.2.0}
without changing a line of its logic. The adapter is responsible only for:
1. unboxing the data frame from the incoming \code{sleep\_diary},
2. calling the unchanged script function,
3. timing the call and diffing rows/columns,
4. calling \code{log\_step()} so the flag ledger stays identical to v1.2.0,
5. boxing the result back into a \code{sleep\_diary}.
\end{Description}
%
\begin{Details}
This is the "wrapper-first" half of the v1.3.0 S3 migration. Once snapshot
tests pin each step's output, the bodies can be internalised one at a time
without the caller ever noticing.
\end{Details}
\HeaderA{plot.bland\_altman}{Plot a Bland-Altman result}{plot.bland.Rul.altman}
%
\begin{Description}
Plot a Bland-Altman result
\end{Description}
%
\begin{Usage}
\begin{verbatim}
## S3 method for class 'bland_altman'
plot(x, ...)
\end{verbatim}
\end{Usage}
%
\begin{Arguments}
\begin{ldescription}
\item[\code{x}] A list returned by \code{bland\_altman()}.

\item[\code{...}] Passed to \code{print.ggplot}.
\end{ldescription}
\end{Arguments}
%
\begin{Value}
The ggplot object, invisibly.
\end{Value}
\HeaderA{plot.sleep\_diary}{Plot the state of a sleep\_diary}{plot.sleep.Rul.diary}
%
\begin{Description}
Degrades gracefully: with ggplot2 available it draws the record count and
flag composition across the recorded chain; without it, falls back to a base
R barplot. Never errors just because a Suggests package is absent.
\end{Description}
%
\begin{Usage}
\begin{verbatim}
## S3 method for class 'sleep_diary'
plot(x, y = NULL, ...)
\end{verbatim}
\end{Usage}
%
\begin{Arguments}
\begin{ldescription}
\item[\code{x}] A \code{sleep\_diary} object.

\item[\code{y}] Unused; present for S3 generic consistency.

\item[\code{...}] Unused.
\end{ldescription}
\end{Arguments}
%
\begin{Value}
The plot object, invisibly.
\end{Value}
\HeaderA{plot.threshold\_validation}{Plot a threshold validation report}{plot.threshold.Rul.validation}
%
\begin{Description}
Draws the Bland-Altman plots for SOL and WASO side by side, overlaid
with the configured threshold lines to show their position relative to
the measurement-noise envelope.
\end{Description}
%
\begin{Usage}
\begin{verbatim}
## S3 method for class 'threshold_validation'
plot(x, ...)
\end{verbatim}
\end{Usage}
%
\begin{Arguments}
\begin{ldescription}
\item[\code{x}] A \code{threshold\_validation} object.

\item[\code{...}] Unused.
\end{ldescription}
\end{Arguments}
%
\begin{Value}
The combined ggplot object, invisibly.
\end{Value}
\HeaderA{print.bland\_altman}{Print a Bland-Altman summary}{print.bland.Rul.altman}
%
\begin{Description}
Print a Bland-Altman summary
\end{Description}
%
\begin{Usage}
\begin{verbatim}
## S3 method for class 'bland_altman'
print(x, ...)
\end{verbatim}
\end{Usage}
%
\begin{Arguments}
\begin{ldescription}
\item[\code{x}] A list returned by \code{bland\_altman()}.

\item[\code{...}] Unused.
\end{ldescription}
\end{Arguments}
%
\begin{Value}
\code{x}, invisibly.
\end{Value}
\HeaderA{print.sleep\_diary}{One-line status of a sleep\_diary}{print.sleep.Rul.diary}
%
\begin{Description}
One-line status of a sleep\_diary
\end{Description}
%
\begin{Usage}
\begin{verbatim}
## S3 method for class 'sleep_diary'
print(x, ...)
\end{verbatim}
\end{Usage}
%
\begin{Arguments}
\begin{ldescription}
\item[\code{x}] A \code{sleep\_diary} object.

\item[\code{...}] Unused.
\end{ldescription}
\end{Arguments}
%
\begin{Value}
\code{x}, invisibly.
\end{Value}
\HeaderA{print.summary.sleep\_diary}{Print a sleep\_diary summary}{print.summary.sleep.Rul.diary}
%
\begin{Description}
Print a sleep\_diary summary
\end{Description}
%
\begin{Usage}
\begin{verbatim}
## S3 method for class 'summary.sleep_diary'
print(x, ...)
\end{verbatim}
\end{Usage}
%
\begin{Arguments}
\begin{ldescription}
\item[\code{x}] A \code{summary.sleep\_diary} object.

\item[\code{...}] Unused.
\end{ldescription}
\end{Arguments}
%
\begin{Value}
\code{x}, invisibly.
\end{Value}
\HeaderA{print.threshold\_validation}{Print a threshold validation report}{print.threshold.Rul.validation}
%
\begin{Description}
Print a threshold validation report
\end{Description}
%
\begin{Usage}
\begin{verbatim}
## S3 method for class 'threshold_validation'
print(x, ...)
\end{verbatim}
\end{Usage}
%
\begin{Arguments}
\begin{ldescription}
\item[\code{x}] A \code{threshold\_validation} object from \code{validate\_thresholds()}.

\item[\code{...}] Unused.
\end{ldescription}
\end{Arguments}
%
\begin{Value}
\code{x}, invisibly.
\end{Value}
\HeaderA{run\_cleaning\_chain}{The S3 cleaning chain}{run.Rul.cleaning.Rul.chain}
%
\begin{Description}
Runs the pure data-in / data-out portion of the pipeline as a single
composable chain. These are the six steps that take a data frame and return
a data frame with no file I/O and no reliance on the global environment:
timestamps, intervals, sequence normalisation, manual corrections, duration
corrections and metric computation.
\end{Description}
%
\begin{Usage}
\begin{verbatim}
run_cleaning_chain(
  data,
  corrections_df = data.frame(),
  manual_unusual_df = data.frame(),
  cfg = NULL,
  verbose = getOption("sleepcleanr.verbose", TRUE)
)
\end{verbatim}
\end{Usage}
%
\begin{Arguments}
\begin{ldescription}
\item[\code{data}] A raw data frame, or an existing \code{sleep\_diary} object.

\item[\code{corrections\_df}] Data frame of manual error corrections (Step 6).
Pass an empty \code{data.frame()} to run without corrections.

\item[\code{manual\_unusual\_df}] Data frame of manual unusual-pattern decisions.

\item[\code{cfg}] List or NULL. Pipeline configuration. Defaults to the
\code{pipeline\_config} published by \code{run\_pipeline()}, if present.

\item[\code{verbose}] Logical. Print per-step progress.
\end{ldescription}
\end{Arguments}
%
\begin{Details}
The remaining steps of \code{run\_pipeline()} -- data loading (1), the
field-misentry check (1.5), review-file generation (5), second-review
consensus (5.75), auto-detection (8), the cross-participant check (8.5) and
visualisation (9) -- are deliberately NOT part of the chain in v1.3.0. They
write files or publish objects into the global environment, so folding them
in would change behaviour rather than just re-shape it. They stay in
\code{run\_pipeline()} until snapshot tests cover them.
\end{Details}
%
\begin{Value}
A \code{sleep\_diary} object carrying the cleaned data and the full
step history. Call \code{as.data.frame()} on it to recover a plain data
frame identical in shape to the v1.2.0 \code{corrected\_ema\_data}.
\end{Value}
%
\begin{Examples}
\begin{ExampleCode}
## Not run: 
cleaned <- run_cleaning_chain(raw_df, corrections, unusual)
summary(cleaned)
plot(cleaned)
corrected_ema_data <- as.data.frame(cleaned)

## End(Not run)
\end{ExampleCode}
\end{Examples}
\HeaderA{run\_figure\_index}{Regenerate the figure\_index.png contact sheet}{run.Rul.figure.Rul.index}
%
\begin{Description}
Regenerate the figure\_index.png contact sheet
\end{Description}
%
\begin{Usage}
\begin{verbatim}
run_figure_index(viz_dir = "latest_visualization")
\end{verbatim}
\end{Usage}
%
\begin{Arguments}
\begin{ldescription}
\item[\code{viz\_dir}] Character. Path to the figure run directory. Default
\code{"latest\_visualization"}.
\end{ldescription}
\end{Arguments}
%
\begin{Value}
Invisibly TRUE.
\end{Value}
\HeaderA{run\_pipeline}{Run the full SPL Sleep pipeline}{run.Rul.pipeline}
%
\begin{Description}
Executes the complete sleep EMA data cleaning pipeline.
Steps 2--7 now flow through the S3 chain (v1.3.1), giving every step
automatic provenance tracking, contract assertions, and a 2.6x speed-up.
Steps that write files or read human-reviewed CSVs remain as direct
\code{source()} calls for backward compatibility.
\end{Description}
%
\begin{Usage}
\begin{verbatim}
run_pipeline(
  config = NULL,
  project_dir = ".",
  skip_visualization = FALSE,
  finalize = TRUE,
  verbose = TRUE
)
\end{verbatim}
\end{Usage}
%
\begin{Arguments}
\begin{ldescription}
\item[\code{config}] Character or list. Path to a YAML config file, or a config
list (from \code{load\_config()}). If NULL, uses the bundled default.

\item[\code{project\_dir}] Character. Path to the project root. Default ".".

\item[\code{skip\_visualization}] Logical. If TRUE, skip visualization.

\item[\code{finalize}] Logical. If TRUE (default) run \code{finalize\_columns()} as
Step 10 and write the delivered datasets. Set FALSE to stop after the
cleaning run and inspect \code{corrected\_ema\_data} yourself. Before v1.4
this step had to be invoked by hand, which meant a plain
\code{run\_pipeline()} produced no Dataset A or B at all.

\item[\code{verbose}] Logical. Print progress. Default TRUE.
\end{ldescription}
\end{Arguments}
%
\begin{Value}
Invisibly returns TRUE on successful completion.
\end{Value}
\HeaderA{run\_report}{Run the reporting stage}{run.Rul.report}
%
\begin{Description}
Run the reporting stage
\end{Description}
%
\begin{Usage}
\begin{verbatim}
run_report(config = NULL, project_dir = ".")
\end{verbatim}
\end{Usage}
%
\begin{Arguments}
\begin{ldescription}
\item[\code{config}] Character or list. Path to a config YAML, a configuration
list (from \code{load\_config()}), or NULL for the bundled default.

\item[\code{project\_dir}] Character. Path to the project root. Default ".".
\end{ldescription}
\end{Arguments}
%
\begin{Value}
Invisibly TRUE.
\end{Value}
\HeaderA{run\_setup}{Run the setup-only stage (package / input-file checks)}{run.Rul.setup}
%
\begin{Description}
Checks R packages and input files without loading or cleaning data.
\end{Description}
%
\begin{Usage}
\begin{verbatim}
run_setup(config = NULL, project_dir = ".")
\end{verbatim}
\end{Usage}
%
\begin{Arguments}
\begin{ldescription}
\item[\code{config}] Character or list. Path to a config YAML, a configuration
list (from \code{load\_config()}), or NULL for the bundled default.

\item[\code{project\_dir}] Character. Path to the project root. Default ".".
\end{ldescription}
\end{Arguments}
%
\begin{Value}
Invisibly TRUE.
\end{Value}
\HeaderA{run\_synthetic\_demo}{Run the complete pipeline on bundled synthetic demo data}{run.Rul.synthetic.Rul.demo}
%
\begin{Description}
Convenience function that runs the full pipeline using the synthetic
EMA diary dataset bundled with sleepcleanr. Useful for testing, demos,
and validation without needing real data.
\end{Description}
%
\begin{Usage}
\begin{verbatim}
run_synthetic_demo(project_dir = ".", verbose = TRUE)
\end{verbatim}
\end{Usage}
%
\begin{Arguments}
\begin{ldescription}
\item[\code{project\_dir}] Character. Path to the project root (where output/ will be created).
Default ".".

\item[\code{verbose}] Logical. Print progress. Default TRUE.
\end{ldescription}
\end{Arguments}
%
\begin{Details}
Synthetic data includes deliberately injected errors across all rule categories
to benchmark detection and correction performance. Results are written to
\code{output/latest\_visualization\_synth\_nXXX/} in the current working directory.
\end{Details}
%
\begin{Value}
Invisibly TRUE on successful completion.
\end{Value}
%
\begin{Examples}
\begin{ExampleCode}
## Not run: 
  # Run the full pipeline on synthetic data in the current directory
  run_synthetic_demo()

  # Run in a specific directory
  run_synthetic_demo(project_dir = "~/my_sleepcleanr_run")

## End(Not run)
\end{ExampleCode}
\end{Examples}
\HeaderA{run\_visualization}{Run only the visualization stage on already-cleaned data}{run.Rul.visualization}
%
\begin{Description}
Loads config + inputs and regenerates the diagnostic figures.
\end{Description}
%
\begin{Usage}
\begin{verbatim}
run_visualization(config = NULL, project_dir = ".")
\end{verbatim}
\end{Usage}
%
\begin{Arguments}
\begin{ldescription}
\item[\code{config}] Character or list. Path to a config YAML, a configuration
list (from \code{load\_config()}), or NULL for the bundled default.

\item[\code{project\_dir}] Character. Path to the project root. Default ".".
\end{ldescription}
\end{Arguments}
%
\begin{Value}
Invisibly TRUE.
\end{Value}
\HeaderA{sleep\_diary}{The sleep\_diary S3 class}{sleep.Rul.diary}
%
\begin{Description}
A thin container that carries the working data frame plus the provenance
of every pipeline step that has touched it.  Each step wraps its
unchanged v1.2.0 logic, only adding timing, column-diff tracking and
ledger logging.
\end{Description}
%
\begin{Details}
Constructors, predicates and validation are documented individually:
\code{\LinkA{new\_sleep\_diary}{new.Rul.sleep.Rul.diary}}, \code{\LinkA{as\_sleep\_diary}{as.Rul.sleep.Rul.diary}},
\code{\LinkA{is\_sleep\_diary}{is.Rul.sleep.Rul.diary}}, \code{\LinkA{validate\_sleep\_diary}{validate.Rul.sleep.Rul.diary}}.
\end{Details}
\HeaderA{STANDARD\_LEVELS}{Fixed category vocabularies for the ledger}{STANDARD.Rul.LEVELS}
%
\begin{Description}
Fixed category vocabularies for the ledger
\end{Description}
%
\begin{Usage}
\begin{verbatim}
STANDARD_LEVELS
\end{verbatim}
\end{Usage}
\HeaderA{step\_apply\_corrections}{Step 6 -- apply manual corrections}{step.Rul.apply.Rul.corrections}
%
\begin{Description}
Step 6 -- apply manual corrections
\end{Description}
%
\begin{Usage}
\begin{verbatim}
step_apply_corrections(x, corrections_df, manual_unusual_df)
\end{verbatim}
\end{Usage}
%
\begin{Arguments}
\begin{ldescription}
\item[\code{x}] A \code{sleep\_diary} object.

\item[\code{corrections\_df}] Data frame of manual error corrections.

\item[\code{manual\_unusual\_df}] Data frame of manual unusual-pattern decisions.
\end{ldescription}
\end{Arguments}
%
\begin{Value}
A \code{sleep\_diary} object.
\end{Value}
\HeaderA{step\_apply\_duration\_corrections}{Step 6.5 -- apply duration corrections}{step.Rul.apply.Rul.duration.Rul.corrections}
%
\begin{Description}
Chains the three v1.2.0 correction appliers in their original order:
nap/exercise, then sleep-metric durations, then human metric acceptances.
\end{Description}
%
\begin{Usage}
\begin{verbatim}
step_apply_duration_corrections(x)
\end{verbatim}
\end{Usage}
%
\begin{Arguments}
\begin{ldescription}
\item[\code{x}] A \code{sleep\_diary} object.
\end{ldescription}
\end{Arguments}
%
\begin{Value}
A \code{sleep\_diary} object.
\end{Value}
\HeaderA{step\_compute\_metrics}{Step 7 -- compute derived sleep metrics}{step.Rul.compute.Rul.metrics}
%
\begin{Description}
Produces the four public contract columns introduced in v1.2.0
(\code{sleep\_efficiency\_pct}, \code{sol\_h}, \code{waso\_h},
\code{sleep\_duration\_h}) that the flag evaluators consume.
\end{Description}
%
\begin{Usage}
\begin{verbatim}
step_compute_metrics(x)
\end{verbatim}
\end{Usage}
%
\begin{Arguments}
\begin{ldescription}
\item[\code{x}] A \code{sleep\_diary} object.
\end{ldescription}
\end{Arguments}
%
\begin{Value}
A \code{sleep\_diary} object.
\end{Value}
\HeaderA{step\_normalize\_sequence}{Step 4 -- normalise sleep time sequence}{step.Rul.normalize.Rul.sequence}
%
\begin{Description}
Step 4 -- normalise sleep time sequence
\end{Description}
%
\begin{Usage}
\begin{verbatim}
step_normalize_sequence(x, flip_gap_hours = NULL, swap_threshold_hours = NULL)
\end{verbatim}
\end{Usage}
%
\begin{Arguments}
\begin{ldescription}
\item[\code{x}] A \code{sleep\_diary} object.

\item[\code{flip\_gap\_hours}] Numeric. Gap above which an AM/PM flip is applied.
Defaults to the configured \code{timestamp.sequence.max\_gap\_hours}, or 12.

\item[\code{swap\_threshold\_hours}] Numeric. Absolute gap below which a minor
order error is corrected by swapping the pair. Defaults to the
configured \code{normalize.swap\_threshold\_hours}, or 3.
\end{ldescription}
\end{Arguments}
%
\begin{Value}
A \code{sleep\_diary} object.
\end{Value}
\HeaderA{step\_process\_intervals}{Step 3 -- parse interval durations}{step.Rul.process.Rul.intervals}
%
\begin{Description}
Step 3 -- parse interval durations
\end{Description}
%
\begin{Usage}
\begin{verbatim}
step_process_intervals(
  x,
  vars = c("duration_totalmin_sol_estimate_am", "duration_totalmin_waso_estimate_am",
    "duration_totalmin_napstoday_PM", "exercisetoday_PM_totalmin_Light",
    "exercisetoday_PM_totalmin_Moderate", "exercisetoday_PM_totalmin_Vigorous",
    "exercisetoday_PM_totalmin_Strength")
)
\end{verbatim}
\end{Usage}
%
\begin{Arguments}
\begin{ldescription}
\item[\code{x}] A \code{sleep\_diary} object.

\item[\code{vars}] Character vector of interval variables to process.
\end{ldescription}
\end{Arguments}
%
\begin{Value}
A \code{sleep\_diary} object.
\end{Value}
\HeaderA{step\_process\_timestamps}{Step 2 -- parse timestamps}{step.Rul.process.Rul.timestamps}
%
\begin{Description}
Step 2 -- parse timestamps
\end{Description}
%
\begin{Usage}
\begin{verbatim}
step_process_timestamps(
  x,
  vars = c("time_bed_am", "time_sleep_am", "time_awake_am", "time_getup_am",
    "caffeinetoday_PM", "alcoholtoday_PM", "nicotine_amount_pm", "cannabis_amount_pm")
)
\end{verbatim}
\end{Usage}
%
\begin{Arguments}
\begin{ldescription}
\item[\code{x}] A \code{sleep\_diary} object.

\item[\code{vars}] Character vector of timestamp variables to process.
\end{ldescription}
\end{Arguments}
%
\begin{Value}
A \code{sleep\_diary} object.
\end{Value}
\HeaderA{summarise\_missing}{Summarise missing-data patterns per participant}{summarise.Rul.missing}
%
\begin{Description}
Summarise missing-data patterns per participant
\end{Description}
%
\begin{Usage}
\begin{verbatim}
summarise_missing(data, group_col = "pid")
\end{verbatim}
\end{Usage}
%
\begin{Arguments}
\begin{ldescription}
\item[\code{data}] A data frame after calling \code{handle\_missing()}.

\item[\code{group\_col}] Character. Participant column.
\end{ldescription}
\end{Arguments}
%
\begin{Value}
A data frame: per participant, counts of each missing reason
and the number of LOCF-filled values.
\end{Value}
\HeaderA{summarise\_outliers}{Summarise IQR outlier flags}{summarise.Rul.outliers}
%
\begin{Description}
Summarise IQR outlier flags
\end{Description}
%
\begin{Usage}
\begin{verbatim}
summarise_outliers(data, group_col = "pid")
\end{verbatim}
\end{Usage}
%
\begin{Arguments}
\begin{ldescription}
\item[\code{data}] A data frame after calling \code{flag\_statistical\_outliers()}.

\item[\code{group\_col}] Character. Participant column.
\end{ldescription}
\end{Arguments}
%
\begin{Value}
A data frame: one row per participant with flag counts per metric.
\end{Value}
\HeaderA{summary.sleep\_diary}{Tabulate the whole pipeline chain recorded in a sleep\_diary}{summary.sleep.Rul.diary}
%
\begin{Description}
Returns one row per step: how many records went in and out, how many columns
the step added, and how long it took. When the flag ledger from
\code{log\_step()} is populated it is joined on, so the same call answers both
"what did each step do" and "how many records were flagged after it".
\end{Description}
%
\begin{Usage}
\begin{verbatim}
## S3 method for class 'sleep_diary'
summary(object, ...)
\end{verbatim}
\end{Usage}
%
\begin{Arguments}
\begin{ldescription}
\item[\code{object}] A \code{sleep\_diary} object.

\item[\code{...}] Unused.
\end{ldescription}
\end{Arguments}
%
\begin{Value}
A data frame with one row per recorded step, invisibly printed.
\end{Value}
\HeaderA{tally\_standard}{Tally one standard's labels into a fixed-level count vector. Returns all-NA (named by levels) if the label vector is entirely NA.}{tally.Rul.standard}
%
\begin{Description}
Tally one standard's labels into a fixed-level count vector.
Returns all-NA (named by levels) if the label vector is entirely NA.
\end{Description}
%
\begin{Usage}
\begin{verbatim}
tally_standard(labels, levels)
\end{verbatim}
\end{Usage}
%
\begin{Arguments}
\begin{ldescription}
\item[\code{labels}] Character vector of labels to tally.

\item[\code{levels}] Character vector of all possible category levels.
\end{ldescription}
\end{Arguments}
%
\begin{Value}
Named integer vector with counts per level.
\end{Value}
\HeaderA{validate\_columns}{Validate that required columns exist}{validate.Rul.columns}
%
\begin{Description}
Validate that required columns exist
\end{Description}
%
\begin{Usage}
\begin{verbatim}
validate_columns(data, required, label = "data")
\end{verbatim}
\end{Usage}
%
\begin{Arguments}
\begin{ldescription}
\item[\code{data}] Data frame.

\item[\code{required}] Character vector of column names that must exist.

\item[\code{label}] Character. Description of what's being checked (for error message).
\end{ldescription}
\end{Arguments}
%
\begin{Value}
Invisibly TRUE. Stops with error if columns are missing.
\end{Value}
\HeaderA{validate\_column\_types}{Validate column types in a data frame}{validate.Rul.column.Rul.types}
%
\begin{Description}
Checks that specified columns have the expected R types.
\end{Description}
%
\begin{Usage}
\begin{verbatim}
validate_column_types(data, type_spec, label = "data")
\end{verbatim}
\end{Usage}
%
\begin{Arguments}
\begin{ldescription}
\item[\code{data}] Data frame.

\item[\code{type\_spec}] Named list mapping column names to expected types
(e.g. \code{list(pid = "numeric", StartDate = "Date")}).
Use \code{"numeric"}, \code{"character"}, \code{"POSIXct"}, \code{"Date"}.

\item[\code{label}] Character. Description of data being checked.
\end{ldescription}
\end{Arguments}
%
\begin{Value}
Invisibly TRUE. Stops with error on mismatch.
\end{Value}
\HeaderA{validate\_no\_r\_code\_in\_paths}{Validate config file paths for R code expressions}{validate.Rul.no.Rul.r.Rul.code.Rul.in.Rul.paths}
%
\begin{Description}
Checks that data file paths in the config are absolute paths, not
R expressions like paste0(...) or file.path(...). Call during
pipeline setup to catch config errors early.
\end{Description}
%
\begin{Usage}
\begin{verbatim}
validate_no_r_code_in_paths(cfg)
\end{verbatim}
\end{Usage}
%
\begin{Arguments}
\begin{ldescription}
\item[\code{cfg}] Config list (from load\_config)
\end{ldescription}
\end{Arguments}
\HeaderA{validate\_schema}{Canonical input schema validator}{validate.Rul.schema}
%
\begin{Description}
Single source of truth for raw input columns. Call right after
\code{adapt\_columns()} so missing or misnamed columns fail loudly.
\end{Description}
%
\begin{Usage}
\begin{verbatim}
validate_schema(data, config, label = "raw EMA input (post-adaptation)")
\end{verbatim}
\end{Usage}
%
\begin{Arguments}
\begin{ldescription}
\item[\code{data}] A data frame to validate.
\item[\code{config}] Pipeline configuration list.
\item[\code{label}] Character label for error messages.
\end{ldescription}
\end{Arguments}
%
\begin{Details}
Design: each schema entry is a logical field resolved through the config
column mapping. The validator accepts either the mapped internal key or the
raw default name, tolerating the current config-key vs hardcoded-name
mismatch.
\end{Details}
\HeaderA{validate\_sleep\_diary}{Validate a sleep\_diary object}{validate.Rul.sleep.Rul.diary}
%
\begin{Description}
Checks the structural invariants the pipeline relies on. Called by the step
adapters so a malformed object fails loudly at the boundary rather than
silently three steps later.
\end{Description}
%
\begin{Usage}
\begin{verbatim}
validate_sleep_diary(x)
\end{verbatim}
\end{Usage}
%
\begin{Arguments}
\begin{ldescription}
\item[\code{x}] An object to validate.
\end{ldescription}
\end{Arguments}
%
\begin{Value}
\code{x}, invisibly, if valid; otherwise an error is raised.
\end{Value}
\HeaderA{validate\_thresholds}{Validate cleaning thresholds against Bland-Altman agreement limits}{validate.Rul.thresholds}
%
\begin{Description}
Runs Bland-Altman analysis on the SOL and WASO self-report / computed
pairs, then compares each threshold in the pipeline configuration against
the 95 
the two measurement methods).
\end{Description}
%
\begin{Usage}
\begin{verbatim}
validate_thresholds(data, cfg = NULL)
\end{verbatim}
\end{Usage}
%
\begin{Arguments}
\begin{ldescription}
\item[\code{data}] A data frame. Typically \code{corrected\_ema\_data} after
\code{run\_cleaning\_chain()} or \code{run\_pipeline()}.

\item[\code{cfg}] A pipeline configuration list (from \code{load\_config()} or
\code{yaml::read\_yaml()}). If \code{NULL}, reads reasonable defaults
from the current \code{pipeline\_config} global.
\end{ldescription}
\end{Arguments}
%
\begin{Details}
A threshold is considered \strong{safe} (\eqn{\checkmark}{}) when it sits
at least \strong{3\eqn{\times}{}} the typical disagreement away from zero bias.
Below 2\eqn{\times}{} the threshold is inside the normal measurement-noise range
and will produce many false positives. Between 2\eqn{\times}{} and 3\eqn{\times}{} is
borderline and warrants a close look at the per-participant data.

This function is \strong{advisory only}. It does not change any threshold
and its output is never fed back into the pipeline decision tree.
\end{Details}
%
\begin{Value}
A data frame with one row per evaluated threshold. Columns:
\code{threshold\_name}, \code{value}, \code{loa\_half\_width},
\code{ratio}, \code{assessment}. The object also carries a
\code{summary} attribute with free-text interpretation and
\code{bland\_altman} attributes holding the raw BA result lists.
\end{Value}
\HeaderA{verification\_run\_dir}{Resolve the stable, never-wiped verification directory for a run}{verification.Rul.run.Rul.dir}
\keyword{internal}{verification\_run\_dir}
%
\begin{Description}
Verification artifacts (S3-vs-legacy snapshot \code{.rds} pairs, the
Markdown verification report, advisory analyses such as Bland-Altman
plots) must survive the *next* pipeline run. The figure run directory
returned by \code{\LinkA{figure\_run\_dir}{figure.Rul.run.Rul.dir}} does not survive it -- every run
of \code{sleep\_visualization.R} deletes and rebuilds that directory from
scratch.
\end{Description}
%
\begin{Usage}
\begin{verbatim}
verification_run_dir(cfg = NULL, data_tag, n_records = NULL)
\end{verbatim}
\end{Usage}
%
\begin{Arguments}
\begin{ldescription}
\item[\code{cfg}] Optional. Pipeline configuration list (preferred; falls back to
the global environment when omitted).

\item[\code{data\_tag}] Character. "real", "synth", or "unknown".

\item[\code{n\_records}] Numeric or NULL. Row count appended to the directory name.
\end{ldescription}
\end{Arguments}
%
\begin{Details}
Earlier, \code{verification/} was nested *inside* the wiped run directory
and rescued around each wipe with a rename-out/rename-back dance
implemented once, in \code{sleep\_visualization.R}. That single
implementation was the only thing standing between the wipe and data loss;
on 2026-08-11 a wipe ran before the dance existed and permanently deleted
that day's verification report (it was gitignored, so it could not be
recovered from git either -- see the "History note" in
\code{output/verification/real\_n13990/VERIFICATION\_2026-08-10.md}). This
function fixes the root cause instead of guarding the symptom: it returns
a path that is a *sibling* of the run directory, not a child of it, so no
wipe -- current or future, in this script or any other -- can reach it.
No preserve logic is required anywhere, and none can be forgotten.
\end{Details}
%
\begin{Value}
Character. Relative path to the stable verification directory,
e.g. \code{"output/verification/real\_n13990"} or
\code{"verification/synth\_n280"}.
\end{Value}
\HeaderA{write\_step\_ledger}{Persist the ledger to CSV (long form)}{write.Rul.step.Rul.ledger}
%
\begin{Description}
Persist the ledger to CSV (long form)
\end{Description}
%
\begin{Usage}
\begin{verbatim}
write_step_ledger(path = "output/step_flag_ledger.csv")
\end{verbatim}
\end{Usage}
%
\begin{Arguments}
\begin{ldescription}
\item[\code{path}] Character. Output CSV file path.
\end{ldescription}
\end{Arguments}
\printindex{}
\end{document}
